\input{../preamble.tex}
\usepackage[normalem]{ulem} % pour sout, texte barré

\SetDate[05/05/2026]

\begin{document}
\pagestyle{fancy}
\fancyhead[L]{Seconde 5}
\fancyhead[C]{\textbf{Statistiques}}
\fancyhead[R]{\today}

\subsection*{Automatismes}

\vspace{-10pt}

% calcul de moyenne


\exe{}{
	Calculer la moyenne de chaque série statistique suivante.
		\begin{multicols}{2}
		\begin{enumerate}
			\item $(0; 20)$
			\item $(0; 0; 20)$
			\item $(0; 20; 20)$
			\item $(-10; 23 ; -10 ; -23 ; 10 ; 10)$
			\item $(-50 ; -100 ; -23 ; -31; -50 ; -50)$
			\item $(0; 1;2;3; \dots; 18; 19; 20)$
		\end{enumerate}
		\end{multicols}

}{exe:0}{

		\begin{enumerate}
			\item $\overline{X} = \dfrac{0+20}{2} = 10$
			\item $\overline{X} = \dfrac{0+0+20}{3} = \dfrac{20}3 \approx 6,67$
			\item $\overline{X} = \dfrac{0+20+20}{3} = \dfrac{40}3 \approx 13,33$
			\item $\overline{X} = \dfrac{-10+23-10-23+10+10}{6} = 0$. La série est symétrique par rapport à $0$ (chaque valeur a son opposé du même effectif).
			\item $\overline{X} = \dfrac{-50-100-23-31-50-50}{6} = -\dfrac{304}6 = -\dfrac{152}3 \approx -50,67$
			\item $\overline{X} = \dfrac{0+ 1+2+3+ \dots+ 18+ 19+ 20}{21} = 10 $. La série est symétrique par rapport à $10$.
		\end{enumerate}
		
		Cette dernière série se généralise pour donner une formule compacte des \emph{nombres triangulaires}. 
		C'est l'objet de l'exercice \ref{exe:somme-triangle} d'appronfondissement.
}

\vspace{-20pt}

\subsection*{Exercices}

\exe{}{
	Assil décide gracieusement de donner un point de sa note à super copine Lina. Est-ce que la moyenne de la classe est toujours la même ?
}{exe:assil-lina}{
	Oui car le nombre de notes et la somme des notes restent inchangés !
}


\exe{}{\label{ex:33}
	Pour chaque série statistique suivante, calculer la moyenne.
		\begin{multicols}{2}
		\begin{enumerate}
			\item 
				\begin{tabular}{|c|c|c|c|c|}\hline
				Valeur   & 9 & 13 & 11 & 7 \\ \hline
				Effectif & 5 & 12 & 5 & 12 \\ \hline
				\end{tabular}
				
			\item 
				\begin{tabular}{|c|c|c|c|c|}\hline
				Valeur   & 2 & 13 & 18 & 7 \\ \hline
				Effectif & 5 & 12 & 5 & 12 \\ \hline
				\end{tabular}
				
			\item 
				\begin{tabular}{|c|c|c|c|c|}\hline
				Valeur   & 0 & 20 & 2 & 5 \\ \hline
				Effectif & 5 & 14 & 10 & 2  \\ \hline
				\end{tabular}

			\item 
				\begin{tabular}{|c|c|c|c|}\hline
				Valeur   & 0 & 20 & 10 \\ \hline
				Effectif & 17 & 17 & 1 \\ \hline
				\end{tabular}


		\end{enumerate}
		\end{multicols}

}{exe:1}{
		\begin{enumerate}
			\item 
				$\overline{X} = \dfrac{9\cdot5 + 13\cdot12 + 11\cdot5 + 7\cdot12}{5+12+5+12} = 10$
				
			\item 
				$\overline{X} = \dfrac{2\cdot5 + 13\cdot12 + 18\cdot5 + 7\cdot12}{5+12+5+12} = 10$
				
			\item 
				$\overline{X} = \dfrac{0\cdot5 + 20\cdot14 + 2\cdot10 + 5\cdot2}{5+14+10+2} = 10$

			\item 
				$\overline{X} = \dfrac{0\cdot17 + 20\cdot17 + 10\cdot1}{17+17+1} = 10$

		\end{enumerate}

}

\exemulticols{, difficulty=1}{
	Pour quel $N\in\N$ la moyenne de la série ci-contre est-elle 10 ?
}{
		\begin{center}
		\begin{tabular}{|c|c|c|c|c|}\hline
		Valeur   & 0 & 16 & 8 & 12 \\ \hline
		Effectif & 2 & 5 & N & 7  \\ \hline
		\end{tabular}
		\end{center}
}{exe:08}{
	Il existe deux façons de procéder : la recherche locale ou la résolution d'une équation.
	
	\underline{Première méthode}
	
	Explorons d'abord la première méthode. Il s'agit de tester une première valeur de $N$, et de décider de la direction (augmenter $N$ ou diminuer $N$) qui permettra de s'approcher de la solution.
	
	Si $N = 0$, alors la moyenne est donnée par
		\[ \frac{0\times2+16\times5+8\times0+12\times7}{2+5+0+7} \approx 11,7. \]
	Comme le but étant d'atteindre 10, il est cohérent d'augmenter le nombre de notes en-dessous de 10 pour baisser la moyenne.
	En augmentant $N$, au augmente donc l'effectif de la note 8, qui va baisser la moyenne.
	
	Si $N=10$, alors la moyenne est donnée par
		\[ \frac{0\times2+16\times5+8\times10+12\times7}{2+5+10+7} \approx 10,2. \]
	Ceci ne suffit pas mais nous évoluons dans la bonne direction !
	
	Si $N=20$, alors la moyenne est donnée par
		\[ \frac{0\times2+16\times5+8\times20+12\times7}{2+5+20+7} \approx 9,5. \]
	Nous sommes allés trop loin : il faut diminuer $N$ pour augmenter la moyenne.
	
	Comme $N \in ]10 ; 20[$, on peut tester $N=15$ et décider d'une direction pour diviser le nombre de valeurs candidates par deux : c'est la dichotomie.

	Si $N=15$, alors la moyenne est donnée par
		\[ \frac{0\times2+16\times5+8\times15+12\times7}{2+5+15+7} \approx 9,79. \]
	Par conséquent, $N \in ]10 ; 15[$.
	
	Si $N=12$, alors la moyenne est donnée par
		\[ \frac{0\times2+16\times5+8\times12+12\times7}{2+5+12+7} = 10, \]
	et nous avons trouvés notre $N$ !

	\underline{Deuxième méthode}

	La contrainte $\overline{X} = 10$ se traduit par l'égalité suivante à résoudre pour $N\in\N$.
	
		\begin{align*}
			\overline{X} &= 10 \\
			\frac{0\times2+16\times5+8\times N+12\times7}{2+5+N+7} &= 10 \\
			0  + 80 + 8\cdot N + 84 &= 10 \cdot(14 + N) \\
			164 + 8 \cdot N &= 140 + 10 \cdot N \\
			164-140 &= 2 \cdot N \\
			N &= 12		
		\end{align*}

	$12$ est bien un entier naturel, donc $N=12$ est l'unique solution.
	Remarquons d'ailleurs que nousa vons démontré l'unicité de la solution en plus de l'avoir obtenue.
}

%\exe{, difficulty=2}{
%	Est-il possible que la série suivante ait une moyenne de $10$ pour un certain $N\in\N$ ? Si oui, lequel ?
%	
%		\begin{center}
%		\begin{tabular}{|c|c|c|c|c|}\hline
%		Valeur   & 23 & 12 & 2 & 7 \\ \hline
%		Effectif & 5 & 14 & N & 2  \\ \hline
%		\end{tabular}
%		\end{center}
%}{exe:}{
%	La contrainte $\overline{X} = 10$ se traduit par l'égalité suivante à résoudre pour $N\in\N$.
%	
%		\begin{align*}
%			\overline{X} &= 10 \\
%			\dfrac{5\cdot23 + 14\cdot12 + N\cdot2 + 2\cdot7}{5+14+N+2} &= 10 \\
%			115  + 168 + 2\cdot N + 14 &= 10 \cdot(21 + N) \\
%			297 + 2 \cdot N &= 210 + 10 \cdot N \\
%			87 &= 8 \cdot N \\
%			N &= \dfrac{87}8		
%		\end{align*}
%
%	L'unique solution candidate est donc $\dfrac{87}{8}$ qui n'est pas un entier naturel.
%	Donc il n'est pas possible que cette série statistique ait une moyenne de $10$, peu importe l'effectif $N\in\N$.
%	
%	En prenant une valeur approchée, par exemple $N=11$, la moyenne est de $9,96$, pas exactement $10$ !
%}


%
%\exe{}{
%	Pour chaque série statistique suivante, remplir la ligne \og Fréquence \fg \, et calculer la moyenne.
%		\begin{multicols}{2}
%		\begin{enumerate}
%			\item 
%				\begin{tabular}{|c|c|c|}\hline
%				Valeur   & 0 & 20 \\ \hline
%				Effectif & 1 & 2  \\ \hline
%				Fréquence & &  \\ \hline
%				\end{tabular}
%				
%			\item 
%				\begin{tabular}{|c|c|c|}\hline
%				Valeur   & 0 & 20 \\ \hline
%				Effectif & 2 & 1  \\ \hline
%				Fréquence & &  \\ \hline
%				\end{tabular}
%				
%			\item 
%				\begin{tabular}{|c|c|c|c|c|}\hline
%				Valeur   & 11 & 13 & 9 & 7 \\ \hline
%				Effectif & 2 & 1 & 5 & 10 \\ \hline
%				Fréquence & &&&  \\ \hline
%				\end{tabular}
%				
%			\item 
%				\begin{tabular}{|c|c|c|c|c|}\hline
%				Valeur   & -23 & -1 & 1 & 23 \\ \hline
%				Effectif & 67 & 12 & 4 & 17 \\ \hline
%				Fréquence & &&&  \\ \hline
%				\end{tabular}
%
%		\end{enumerate}
%		\end{multicols}
%
%}{exe:2}{
%		\begin{enumerate}
%			\item 
%				La somme des effectifs $N$ est donnée par $N = 1+2 = 3$.
%				Pour trouver la fréquence d'une valeur, on divise son effectif par $N$.
%				
%				\begin{tabular}{|c|c|c|}\hline
%				Valeur   & 0 & 20 \\ \hline
%				Effectif & 1 & 2  \\ \hline
%				Fréquence & $\frac13$ & $\frac23$  \\ \hline
%				\end{tabular}
%				
%				La moyenne est la somme des produits Valeur $\times$ Fréquence.
%				D'où $\overline{X} = 0 \cdot \frac13 + 20 \cdot \frac23 =\dfrac{40}3 \approx 13,33$, qu'on comparera à l'exercice $1$.
%				
%			\item 
%				La somme des effectifs $N$ est donnée par $N = 1+2 = 3$.
%				Pour trouver la fréquence d'une valeur, on divise son effectif par $N$.
%				
%				\begin{tabular}{|c|c|c|}\hline
%				Valeur   & 0 & 20 \\ \hline
%				Effectif & 2 & 1  \\ \hline
%				Fréquence & $\frac23$ & $\frac13$ \\ \hline
%				\end{tabular}
%				
%				D'où $\overline{X} = 0 \cdot \frac23 + 20 \cdot \frac13 =\dfrac{20}3 \approx 6,67$, qu'on comparera à l'exercice $1$.
%				
%			\item 
%				La somme des effectifs $N$ est donnée par $N = 2+1+5+10 = 18$.
%				Pour trouver la fréquence d'une valeur, on divise son effectif par $N$.
%				
%				\begin{tabular}{|c|c|c|c|c|}\hline
%				Valeur   & 11 & 13 & 9 & 7 \\ \hline
%				Effectif & 2 & 1 & 5 & 10 \\ \hline
%				Fréquence & $\frac2{18}$ & $\frac1{18}$ & $\frac5{18}$ & $\frac{10}{18}$  \\ \hline
%				\end{tabular}
%				
%				D'où $\overline{X} = 11\cdot\frac2{18} + 13\cdot \frac1{18} + 9\cdot \frac5{18} 7\cdot\frac{10}{18} = \dfrac{150}{18} \approx 8,33 $.
%				
%			\item 
%				La somme des effectifs $N$ est donnée par $N = 67+12+4+17 = 100$.
%				Pour trouver la fréquence d'une valeur, on divise son effectif par $N$.
%				
%				\begin{tabular}{|c|c|c|c|c|}\hline
%				Valeur   & -23 & -1 & 1 & 23 \\ \hline
%				Effectif & 67 & 12 & 4 & 17 \\ \hline
%				Fréquence & $\frac{67}{100}$ & $\frac{12}{100}$ & $\frac{4}{100}$ & $\frac{17}{100}$ \\ \hline
%				\end{tabular}
%				
%				D'où $\overline{X} = (-23) \cdot \frac{67}{100} + (-1)\cdot \frac{12}{100} + 1\cdot \frac{4}{100} + 23\cdot \frac{17}{100} = -11,58$.
%
%		\end{enumerate}
%
%
%}

\exe{}{
	Un élève reçoit cinq notes listées ainsi : $(4,25 ; 10; 18,5 ; 15,5 ; 11,25)$.
	À ces notes sont associés cinq coefficients, donnés dans le même ordre par $(0,5 ; 1,4 ; 0,1 ; 0,5; 1,1)$.
	\begin{enumerate}
		\item Calculer la moyenne pondérée des notes.
		\item Multiplier tous les coefficients par 10 puis recalculer la moyenne pondérée des notes. Que remarque-t-on ?
	\end{enumerate}
}{exe:3}{
%	Il y a deux manières de faire.
%	Calculons la somme des produits Note $\times$ Coefficient puis diviser par la somme des coefficients.
%		\[ \dfrac{4,25\cdot0,5 + 10\cdot1,4 + 18,5\cdot0,1 + 15,5\cdot 0,5 + 11,25\cdot1,1}{0,5+1,4+0,1+0,5+1,1} = \dfrac{38,1}{3,6}  \approx 10,58. \]
%	D'autre part, on peut calculer le poids de chaque note en normalisant les coefficients, c'est-à-dire en divisant chaque coefficient par leur somme.
%	
%	$p_1 = \dfrac{0,5}{3,6}$, $p_2 = \dfrac{1,4}{3,6}$, $p_3 = \dfrac{0,1}{3,6}$, $p_4 = \dfrac{0,5}{3,6}$, et $p_5 = \dfrac{1,1}{3,6}$, puis
%		\[ 4,25 \cdot p_1 + 10\cdot p_2 + 18,5 \cdot p_3 + 15,5 \cdot p_4 + 11,25\cdot p_5 = 10,58. \]
%
%	La première méthode est plus facile à employer mais elle découle en fait de la seconde : une moyenne est toujours une combinaison convexe de valeurs car les poids somment toujours à $1$.
	\begin{enumerate}
		\item 
		Calculons la somme des produits Note $\times$ Coefficient puis diviser par la somme des coefficients.
			\[ \dfrac{4,25\cdot0,5 + 10\cdot1,4 + 18,5\cdot0,1 + 15,5\cdot 0,5 + 11,25\cdot1,1}{0,5+1,4+0,1+0,5+1,1} = \dfrac{38,1}{3,6}  \approx 10,58. \]
		\item 
		Les nouveaux coefficients deviennent $(5 ; 14 ; 1 ; 5; 11)$ et la nouvelle moyenne est
			\[ \dfrac{4,25\cdot5 + 10\cdot14 + 18,5\cdot1 + 15,5\cdot 5 + 11,25\cdot11}{5+14+1+5+11} = \dfrac{381}{36}  \approx 10,58. \]
		La même moyenne est obtenue, ce qui n'est pas surprenant car le deuxième calcul est obtenu en multipliant le numérateur et le dénominateur du premier calcul par 10, ce qui ne change pas la valeur de la fraction.
	\end{enumerate}
	
	\underline{Conclusion}
	
	Les coefficients donnent un poids à chaque note qui est relatif aux autres coefficients.
	Ainsi, une note coefficient 1 a donc huit fois plus de poids qu'une note coefficient 0,125.
	De la même manière, une note coefficient 8 aura huit fois plus de poids qu'une note coefficient 1.
	
}

% TODO: fix leur/leurs ça me saoule ça je dois apprendre
\exe{,difficulty=1}{
	Un élève reçoit deux notes, 7 et 14, et son méchant professeur souhaite modifier leur coefficient afin que l'élève ait une moyenne inférieure à 10.
	\begin{enumerate}
		\item
		Calculer la moyenne des notes lorsque leurs coefficients sont égaux.
		\item
		Que dire sur les coefficients pour que la moyenne pondérée soit inférieure à 10 ?
		\item
		Trouver deux coefficients tels que la moyenne soit \emph{exactement} de 10.
		
		\emph{Indication : supposer que la somme des coefficients est égale à 1.}
		% est-ce que ça va vraiment aider, ça ?
	\end{enumerate}
}{exe:moyenne-ponderee}{
	\begin{enumerate}
		\item
		D'après l'exercice \ref{exe:3}, lorsque les coefficients sont égaux, ils peuvent être supposés tous les deux égaux à 1.
		On obtient donc une moyenne classique $\frac{7+14}{2} = 10,5$.
		\item
		Pour obtenir une moyenne pondérée inférieure à 10, il faut augmenter le poids de la note la plus basse.
		Le coefficient associé à 7 doit donc être supérieur à celui associé à 14.
		\item
		Posons $c>0$ le coefficent associé à 7.
		D'après l'exerce \ref{exe:3}, nous pouvons diviser les coefficients par leur somme sans changer la moyenne pondérée, et donc supposer que la somme des coefficients vaut 1. 
		Le coefficient associé à 14 est donc $1-c$, et la moyenne pondérée vaut
			\[ \frac{7\times c + 14\times(1-c)}{1} = 7c + 14(1-c). \]
		La contrainte de l'énoncé nous permet donc d'obtenir
			\begin{align*}
				7c+14(1-c) &= 10, \\
				7c + 14 - 14c &= 10, \\
				-7c &= -4, \\
				c &= \frac47.
			\end{align*}
		Les coefficients sont donc $\frac47$ et $\frac37$.
		Notons qu'une stratégie de recherche par dichotomie comme à l'exercice \ref{exe:08} n'aurait pas abouti car les nombres $\frac47$ et $\frac37$ ne sont pas décimaux : leur développement décimal est infini.
	\end{enumerate}
}


\exe{, difficulty=1}{
	En 2025, la moyenne des salaires mensuels d'une entreprise est de 1 860€.
	L'entreprise, après une année où l'inflation se mesurait à 6\% décide d'ajouter des primes aux salaires : tous les employés recevront 50€ mensuellement en plus. 
	%Ces $50$€ sont, eux aussi, soumis à l'inflation.
	
	
	L'inflation désigne une augmentation générale des prix. Ainsi, 6\% d'inflation signifie que, en moyenne, les prix augmentent de 6\%.
	Afin de pouvoir comparer les pouvoirs d'achats \emph{en euros constants},
	il est possible de corriger les salaires de l'inflation en leur appliquant l'évolution réciproque.
	%il est possible de fixer les prix et voir l'inflation plutôt comme une diminution du salaire.
	
	\begin{enumerate}
		\item
		Quel est le coefficient multiplicateur associé à une augmentation de 6\% ?
		\item
		Quel est le coefficient multiplicateur associé à l'évolution réciproque ?
		\item
		Calculer la moyenne des salaires mensuels de l'entreprise en 2026 en euros constants. La prime a-t-elle compensé l'inflation ?
	\end{enumerate}

}{exe:4}{

	\begin{enumerate}
		\item
			\[ CM = 1 + \frac{P}{100} = 1+\frac6{100} = 1,06. \]
		\item
		D'après le cours, c'est $1,06^{-1} = \frac1{1,06}$ et non pas $0,94$ (qui correspond à une diminution de $6\%$) !
		\item
		Notons $X$ la série statistique des salaires avant inflation.
		On utilise ici la linéarité de la moyenne.
		
		La prime s'ajoute après l'inflation, donc la moyenne après inflation et prime s'élève à $\overline{X+50} = \overline{X} + 50$.
		
		Enfin, afin de corriger l'inflation et d'après le texte, on multiplie tous les nouveaux salaires par $(1,06)^{-1}$.
		La nouvelle moyenne est donc donnée par 
			\[ \overline{\frac{1}{1,06} \cdot \left( X + 50 \right) } = \frac{1}{1,06} \cdot \left( \overline{X} + 50 \right) = \frac{1}{1,06} \cdot (1860 + 50) \approx 1801,89. \]
		En euros constants, et en moyenne, la prime ne permet pas de compenser l'inflation.
	\end{enumerate}
	
	
	
}


%
%\exe{}{
%	Calculer et comparer les écart types des séries de l'exercice \ref{ex:33}.
%
%}{exe:5}{
%	On calcule d'abord la variance puis l'écart type.
%	D'après le cours, $\Var{X} = \overline{(X - \overline{X})^2}$ et $\sigma(X) = \sqrt{\Var{X}}$.
%	
%	\begin{enumerate}
%		\item 
%			\[ \Var{X} = \dfrac{5\cdot(9-10)^2 + 12 \cdot(13-10)^2 + 5\cdot(11-10)^2 + 12 \cdot(7-10)^2}{5+12+5+12} = \dfrac{226}{34}. \]
%		
%		D'où $\sigma(X) = \sqrt{\dfrac{226}{34}} \approx 2,578$.
%		
%		\item  
%			\[\Var{X} = \dfrac{5\cdot(2-10)^2 + 12 \cdot(13-10)^2 + 5\cdot(18-10)^2 + 12 \cdot(7-10)^2}{5+12+5+12} = \dfrac{856}{34}. \]
%		
%		D'où $\sigma(X) \approx 5,02$.
%		
%		\item  
%			\[\Var{X} = \dfrac{5\cdot(0-10)^2 + 14 \cdot(20-10)^2 + 10\cdot(2-10)^2 + 2 \cdot(5-10)^2}{5+14+10+2} = \dfrac{2590}{31} . \]
%		
%		D'où $\sigma(X) \approx 9,14.$
%		
%		\item  
%			\[\Var{X} = \dfrac{17\cdot(0-10)^2 + 17 \cdot(20-10)^2 + 1\cdot(10-10)^2 }{17+17+1} = \dfrac{3400}{35}. \]
%		
%		D'où $\sigma(X) \approx 9,86$.
%	\end{enumerate}
%
%	Les écarts types augmentent car les séries statistiques sont de moins en moins concentrées autour de leur moyenne $10$.
%}



%\exe{}{
%	\begin{multicols}{2}
%	On considère la série ci-contre tirée de l'exercice \ref{ex:33}.
%	
%		\begin{center}
%		\begin{tabular}{|c|c|c|c|}\hline
%		Valeur   & 0 & 20 & 10 \\ \hline
%		Effectif & 17 & 17 & 1 \\ \hline
%		\end{tabular}
%		\end{center}
%	\end{multicols}
%		
%	Modifier l'effectif de la valeur $10$ en $10, 20, 50, 100$, et calculer la moyenne et l'écart type de la série obtenue.
%	Comment évolue $\sigma$ ? La série devient-elle plus ou moins concentrée autour de sa moyenne ?
%
%}{exe:6}{
%	Remplaçons l'effectif de $10$ par $k$ pour obtenir directement une formule générale qu'on pourra évaluer en $k=10, 20, 50, 100$.
%
%	D'abord, la moyenne ne change pas :
%		\[ \overline{X} = \dfrac{0\cdot17 + 20\cdot17 + 10\cdot k}{17+17+k} = \dfrac{10\cdot(34 + k)}{34+k} = 10, \]
%	ce qui n'est pas surprenant car ajouter une valeur moyenne ne change pas la moyenne générale.
%
%	On a ainsi 
%		\[ \Var{X} = \dfrac{17\cdot(0-10)^2+17\cdot(20-10)^2 + k\cdot(10-10)^2}{17+17+k} = \dfrac{3400}{34+k}. \]
%		
%	Pour $k=1$, on retrouve l'écart type $\sigma(X) = \sqrt{\dfrac{3400}{35}} \approx 9,86$ de l'exercice précédent.
%	Les autres valeurs sont listées dans le tableau ci-dessous.
%	
%	\begin{center}
%	\begin{tabular}{|c|c|c|c|c|c|}\hline
%		$k$ & 1 & 10 & 20 & 50 & 100 \\ \hline
%		$\sigma(X) \approx$ & 9,86 & 8,79 & 7,93 & 6,36 & 5,04 \\ \hline
%	\end{tabular}
%	\end{center}
%	
%	On remarque que $\sigma$ diminue quand $k$ augmente.
%	En effet, plus le dénominateur de la fraction $\Var{X} = \frac{3400}{34+k}$ augmente, plus la variance diminue, et donc plus l'écart type diminue aussi.
%	
%	La série statistique devient de plus en plus concentrée autour de sa moyenne car il y a de plus en plus de valeurs égales à celle-ci.
%}


% Hyperbole 2nde + ref à Lucile :)
\exe{}{
	Chaque jour du mois d'avril 2026, le \sout{meilleur ami} super cousin de Lucile relève la distance qu'il a parcourue lors de son footing.
	\begin{center}
\def\arraystretch{1}
\setlength\tabcolsep{5pt}
	\begin{tabular}{|*{8}{c|}}\hline
		Distance (km) & 3 & 4 & 5 & 9 & 10 & 12 & 13
		\\\hline
		Effectif & 8 & 6 & 2 & 4 & 3 & 5 & 2
		\\\hline
	\end{tabular}
	\end{center}
	
	\begin{enumerate}
		\item
		Vérifier que la moyenne de cette série est égale à 7 km.
		\item
		On a représenté ci-dessous la série ainsi que les écarts entre les valeurs et la moyenne.
		
		\begin{center}
		\begin{tikzpicture}
		\begin{axis}[
			xmin=0, xmax=14,
			ymin=0, ymax=8,
			x=20pt,
			y=10pt,
			grid=none,
			xtick distance=1,
			axis line style = -,
			extra x ticks={0},
			extra y ticks={0},
			xlabel=Distance (km),
			xlabel near ticks,
			ylabel=Effectif,
			ylabel near ticks,
			ybar,
			bar width=8pt,
		]     
	      \addplot+[bar shift=0,draw=BLACK, thick, fill=BLUE_E!50] plot coordinates {
	        (3,8)
	        (4,6)
	        (5,2)
	        (9,4)
	        (10,3)
	        (12,5)
	        (13,2)
	      };
	      \addplot+[bar shift=0,draw=BLACK, thick, fill=RED_E!50] plot coordinates {
	        (7,8)
	      };
	      
			\draw[GREY_E, stealth-stealth, shorten <=5pt, shorten >=5pt, dashed, yshift=-2pt] (axis cs:3,8) -- (axis cs:7,8);
			\draw[GREY_E, stealth-stealth, shorten <=5pt, shorten >=5pt, dashed, yshift=-2pt] (axis cs:4,6) -- (axis cs:7,6);
			\draw[GREY_E, stealth-stealth, shorten <=5pt, shorten >=5pt, dashed, yshift=-2pt] (axis cs:5,2) -- (axis cs:7,2);
			\draw[GREY_E, stealth-stealth, shorten <=5pt, shorten >=5pt, dashed, yshift=-2pt] (axis cs:7,4) -- (axis cs:9,4);
			\draw[GREY_E, stealth-stealth, shorten <=5pt, shorten >=5pt, dashed, yshift=-2pt] (axis cs:7,3) -- (axis cs:10,3);
			\draw[GREY_E, stealth-stealth, shorten <=5pt, shorten >=5pt, dashed, yshift=-2pt] (axis cs:7,5) -- (axis cs:12,5);
			\draw[GREY_E, stealth-stealth, shorten <=5pt, shorten >=5pt, dashed, yshift=-2pt] (axis cs:7,2) -- (axis cs:13,2);

		\end{axis}
		\end{tikzpicture}
		\end{center}
		\begin{center}
\def\arraystretch{1.5}
\setlength\tabcolsep{15pt}
		\begin{tabular}{|*{8}{c|}}\hline
			Distance (km) & 3 & 4 & 5 & 9 & 10 & 12 & 13
			\\\hline
			Écart & & & &&&&
			\\\hline
			Carré de l'écart & & & & &&&
			\\\hline
			Effectif & 8 & 6 & 2 & 4 & 3 & 5 & 2
			\\\hline
		\end{tabular}
		\end{center}
		
		\begin{enumerate}
			\item
			Compléter la ligne « Écart » avec les écarts de chaque distance à la distance moyenne.
			\item
			% en fait les écarts sont des valeurs absolues donc bon ?
			%Pour n'avoir que des nombres positifs, on calcule le carré de chaque écart.
			%Compléter cette ligne du tableau.
			Calculer le carré de chaque écart et compléter la ligne correspondante du tableau.
			\item
			Calculer la moyenne $V$ des carrés des écarts pondérée par les effectifs.
			Arrondir à $10^{-1}$ près.
			Dans quelle unité est exprimée $V$ ?
			\item
			Calculer $\sigma = \sqrt{V}$ (lu « sigma ») en arrondissant à $10^{-1}$ près.
			Dans quelle unité est \mbox{exprimée $\sigma$ ?}
		\end{enumerate}
	\end{enumerate}
	
}{exe:variance}{
	\begin{enumerate}
		\item
		La moyenne est calculée en faisant
			\[ \frac{3\times8+4\times6+5\times2+9\times4+10\times3+12\times5+13\times2}{8+6+2+4+3+5+2} = \frac{210}{30} = 7. \]
		
		\item\,
		
		\begin{center}
\def\arraystretch{1.5}
\setlength\tabcolsep{15pt}
		\begin{tabular}{|*{8}{c|}}\hline
			Distance (km) & 3 & 4 & 5 & 9 & 10 & 12 & 13
			\\\hline
			Écart & 4 & 3 & 2 & 2 & 3 & 5 & 6
			\\\hline
			Carré de l'écart & 16 & 9 & 4 & 4 & 9 & 25 & 36
			\\\hline
			Effectif & 8 & 6 & 2 & 4 & 3 & 5 & 2
			\\\hline
		\end{tabular}
		\end{center}
		\begin{enumerate}[start=3]
			\item
			La moyenne des carrés des écarts est donnée par
				\[ V = \frac{16\times8+9\times6+4\times2+2\times4+9\times3+25\times5+36\times2}{8+6+2+4+3+5+2} = \frac{430}{30} \approx 14,3. \]
			$V$ est exprimée en $\text{km}^2$.
			\item
			$\sigma = \sqrt{V} \approx \sqrt{14,3} \approx 3,8$, qui est exprimée en km !
		\end{enumerate}
	\end{enumerate}
}



%\subsection*{Exercices supplémentaires}



% Déclic 2nde
\exe{}{
	Deux étudiantes ont obtenues les notes suivantes aux trois derniers partiels de mathématiques : $(10 ; 9 ; 11)$ pour la première et $(6 ; 13 ; 11)$ pour la seconde.
	\begin{enumerate}[label=\roman*)]
		\item
		Calculer la moyenne de chaque étudiante.
		Peut-on dire que leur situation est similaire ?
		\item
		Calculer la variance des notes de chaque étudiante.
		Comment interpréter le résultat ?
		\item
		% TODO: rephrasing
		Une troisième étudiante a pour résultats $(4 ; 12 ; 14)$.
		Sans calculs, décider si l'écart-type de ses notes se plus grand ou plus petit que celui des notes des autres étudiantes.
	\end{enumerate}
}{exe:var-comp}{
	\begin{enumerate}[label=\roman*)]
		\item
		Les deux moyennes sont égales à 10.
		On peut utiliser l'exercice \ref{exe:assil-lina} pour connaître la moyenne sans calcul : ajouter 1 à une note et soustraire 1 à une autre ne change pas la moyenne.
		Ainsi la moyenne de $(10 ; 9 ; 11)$ est égale à la moyenne de $(10 ; 10 ; 10)$.
		D'autre part, la moyenne de $(6 ; 13 ; 11)$ est égale à la moyenne de $(7 ; 12 ; 11)$, et de $(8 ; 11 ; 11)$, et de $(9 ; 10 ; 11)$, et enfin de $(10 ; 10 ; 10)$ !
		
		Les moyennes sont les mêmes mais les situations ne sont pas similaires : la première étudiante a des valeurs plus concentrées autour de 10 que la deuxième.
		La deuxième étudiante a à la fois de bonnes et de mauvaises notes qui se compensent.
		La première n'a que des notes proches de 10.
		\item
		Pour la première étudiante, les carrés des écarts à la moyenne sont donnés par $(0 ; 1 ; 1)$, de moyenne $\frac23 \approx 0,6$.
		
		Pour la deuxième étudiante, les carrés des écarts à la moyenne sont donnés par $(16 ; 9 ; 1)$, de moyenne $\frac{26}3 \approx 8,6$.
		
		La variance est plus grande pour la deuxième étudiante car ses notes sont davantages dispersées autour de la moyenne.
		La variance quantifie la dispersion. Une variance nulle correspondant à une série entièrement moyenne.
		\item
		La moyenne de la troisième étudiante vaut aussi 10.
		Ses notes sont encore plus dispersées autour de celle-ci, notamment à cause du 4 dont le carré de l'écart à la moyenne vaut 64 !
		La variance des notes de la troisième étudiante est supérieur. L'écart-type l'est donc aussi car la fonction racine carrée est croissante :
			\[ V_1 \leq V_2 \implies \sqrt{V_1} \leq \sqrt{V_2}. \]
	\end{enumerate}
}

\exe{}{
	Pour chaque histogramme de la figure \ref{fig:hist} (page 3), calculer en arrondissant à $10^{-2}$
		\begin{multicols}{2}
		\begin{enumerate}[label=\roman*)]
			\item l'effectif total ;
			\item la moyenne ;
			\item la médiane ;
			\item le premier quartile ;
			\item le troisième quartile ; et
			\item l'écart interquartile.
		\end{enumerate}
		\end{multicols}
	Sans connaître la valeur exacte de chaque note, on assignera la valeur moyenne à chaque élément d'une classe.
	Par exemple, de l'histogramme \ref{fig:1a}, on lit 5 notes de valeur 11,5.
}{exe:7}{
	Pour l'histogramme \ref{fig:1a},
	\begin{enumerate}[label=\roman*)]
		\item l'effectif total est de 35
		\item la moyenne est de 
			\[ \frac{
				\splitfrac{5,5+7,5+2\times9,5+2\times10,5+5\times11,5+4\times12,5+2\times13,5}
				{+2\times14,5+4\times15,5+3\times16,5+3\times17,5+5\times18,5+1\times19,5}
				}{35} \approx 14,07. 
			\]
			
		\item 
		La médiane est la note qui sépare la série statistique rangée par ordre croissant en deux.
		Ici, $\frac{35}2 = 17,5$, et donc la 18ème note donnera la médiane.
		En comptant, elle se trouve dans la barre de classe $[14;15[$, et la médiane vaut donc 14,5.
		\item 
		Même raisonnement, mais en divisant l'effectif total par 4.
		$\frac{35}4 = 8,75$, donc la 9è note donne le premier quartile.
		En comptant, elle se trouve dans la barre de classe $[11;12[$, et le premier quartile vaut donc 11,5.
		\item 
		Même raisonnement, mais en prenant les trois quart de l'effectif total.
		$\frac34\cdot35 = 26,25$, donc la 27è note donne le troisième quartile.
		En comptant (par la fin !), elle se trouve dans la barre de classe $[17;18[$, et le troisième quartile vaut donc 17,5.
		\item 
		L'écart interquartile est donné par
			\[ Q_3 - Q_1 = 17,5 - 11,5 = 6. \]
	\end{enumerate}
	
	
	
	Pour l'histogramme \ref{fig:1b},
	\begin{enumerate}[label=\roman*)]
		\item l'effectif total est de 27
		\item la moyenne est de 
			\[ \frac{
				\splitfrac{2,5+3,5+2\times5,5+3\times6,5+7,5+8,5+9,5+2\times10,5+11,5+3\times12,5}
				{+2\times13,5+2\times14,5+2\times15,5+2\times16,5 + 17,5+18,5+19,5}
				}{27} \approx 11,39. 
			\]
			
		\item 
		La médiane est la note qui sépare la série statistique rangée par ordre croissant en deux.
		Ici, $\frac{27}2 = 13,5$, et donc la 14ème note donnera la médiane.
		En comptant, elle se trouve dans la barre de classe $[12;13[$, et la médiane vaut donc 12,5.
		\item 
		Même raisonnement, mais en divisant l'effectif total par 4.
		$\frac{27}4 = 6,75$, donc la 7è note donne le premier quartile.
		En comptant, elle se trouve dans la barre de classe $[6;7[$, et le premier quartile vaut donc 6,5.
		\item 
		Même raisonnement, mais en prenant les trois quart de l'effectif total.
		$\frac27\cdot35 = 20,25$, donc la 21è note donne le troisième quartile.
		En comptant, elle se trouve dans la barre de classe $[15;16[$, et le troisième quartile vaut donc 15,5.
		\item 
		L'écart interquartile est donné par
			\[ Q_3 - Q_1 = 15,5 - 6,5 = 9. \]
	\end{enumerate}
}

\begin{figure}[h!]
  \centering
  \begin{subfigure}[b]{.45\textwidth}
    \centering
  \begin{tikzpicture}[scale=1]
    \begin{axis}[
        ymin=0, ymax=5,
        xmin=0, xmax=20,
        minor x tick num = 1,
        %minor y tick num=1,
        axis line style = -,
        xtick distance=2,
        area style,
        xlabel = {Note sur $20$},
        ylabel = {Effectif},
        xlabel near ticks,
        ylabel near ticks,
		extra x ticks={0},
		extra y ticks={0},
		x=10pt,
		ybar,
		bar width=8pt,
      ]
      \addplot+[bar shift=4pt,draw=BLACK, thick, fill=BLUE_E!50] plot coordinates {
        (5,1)
        (7,1)
        (9,2)
        (10,2)
        (11,5)
        (12,4)
        (13,2)
        (14,2)
        (15,4)
        (16,3)
        (17,3)
        (18,5)
        (19,1)
      };
    \end{axis} 
  \end{tikzpicture}
  \caption{Automatismes 15 à 20.}
  \label{fig:1a}
  \end{subfigure}
  \hfill
  % NOT LINE BREAK!!
	\begin{subfigure}[b]{.45\textwidth}
    \centering
  \begin{tikzpicture}[scale=1]
    \begin{axis}[
        ymin=0, ymax=3,
        xmin=0, xmax=20,
        minor x tick num = 1,
        %minor y tick num=1,
        axis line style = -,
        xtick distance=2,
        area style,
        xlabel = {Note sur $20$},
        ylabel = {Effectif},
        xlabel near ticks,
        ylabel near ticks,
		extra x ticks={0},
		extra y ticks={0},
		x=10pt,
		ybar,
		bar width=8pt,
      ]
      \addplot+[bar shift=4pt,draw=BLACK,thick, fill=GREEN_E!50] plot coordinates {
        (2,1)
        (3,1)
        (5,2)
        (6,3)
        (7,1)
        (8,1)
        (9,1)
        (10,2)
        (11,1)
        (12,3)
        (13,2)
        (14,2)
        (15,2)
        (16,2)
        (17,1)
        (18,1)
        (19,1)
      };
    \end{axis} 
  \end{tikzpicture}
  \caption{Évaluation ``Vecteurs''.}
  \label{fig:1b}
  \end{subfigure}
  
  
  \caption{Histogrammes de notes. Les classes sont de la forme $[k;k+1[$ avec $k$ un entier. 
  %La colonne entre $12$ et $13$ compte donc toutes les notes appartenant à $[12;13[$.
  \\
  \underline{Lecture} : 3 notes de l'évaluation ``Vecteurs'' appartiennent à l'intervalle $[12;13[$.
  }
  \label{fig:hist}
\end{figure}

\newpage

\exe{}{
	Pour chaque histogramme de la figure \ref{fig:hist}, faire une phrase décrivant la série statistique en s'appuyant sur les données obtenues à l'exercice \ref{exe:7}. 
	
	Comparer les deux séries statistiques en identifiant ce qu'elles ont en commun \mbox{et ce qui les distingue.}
}{exe:8}{
	La série de l'histogramme \ref{fig:1a} est une série de 35 notes entre 0 et 20 ayant pour moyenne 14,07.
	La médiane est de 14,5 : la moitié des notes sont inférieures à 14,5 et la moitié sont supérieures.
	Le premier quartile est de 11,5, donc un trois quart des notes sont supérieures à 11,5.
	
	La série de l'histogramme \ref{fig:1a} est une série de 27 notes entre 0 et 20 ayant pour moyenne 11,39.
	Le troisième quartile est de 15,5, donc un quart des notes sont supérieures à 15,5, et trois quart sont inférieures à 11,5.
	Le premier quartile est de 6,5, donc la moitié des valeurs de la série se situent entre 6,5 et 15,5.
	
	La première série a un premier quartile beaucoup plus élevé que celui de la deuxième qui admet donc davantage de petites valeurs.
	La première série a un écart interquartile plus bas que celui de la deuxième qui a donc des valeurs centrales plus proches les unes des autres.
}

\vspace{-10pt}

\subsection*{Approfondissements}

% this is great but i'll do i in the statistics way
%\begin{multicols}{2}
%	\begin{Exercise}[difficulty=2, label=exe:somme-triangle]
%		Pour $n\in\N, n \geq 2$, on considère le carré $(n+1) \times (n+1)$ ci-contre, découpé en $(n+1)^2$ carrés unités.
%		En comptant les carrés sur chaque diagonale, démontrer que
%			\[ 1 + 2 + 3 + \dots + (n-1) + n = \dfrac{n(n+1)}2. \]
%		En déduire la valeur de la somme des 10 000 premiers entiers
%			\[ 1 + 2 + 3 + \hdots + 10~000. \]
%	\end{Exercise}
%	\begin{center}
%	\begin{tikzpicture}[scale=0.4]
%		\foreach \r in {0, ..., 10} {
%			\draw[black,thick] (0,\r) -- (10,\r);
%			\draw[black, thick] (\r,0) -- (\r, 10);
%		}
%		
%		\draw[decoration={brace},decorate]
%  			(0,10.2) -- node[above] {$n+1$} (10,10.2);
%		\draw[decoration={brace, mirror},decorate]
%  			(10.2,0) -- node[right] {$n+1$} (10.2,10);
%  			
%		\draw[->] (-0.5, 1.5) -- (1.5,-0.5) node[below] {$1$};
%		\draw[->] (-0.5, 2.5) -- (2.5,-.5) node[below] {$2$};
%		\draw[->] (-0.5, 3.5) -- (3.5,-.5) node[right=12pt, below] {$3$ $\cdots$};
%	
%	\end{tikzpicture}
%	\end{center}
%\end{multicols}
%\begin{Answer}[ref=exe:somme-triangle]
%	todo
%\end{Answer}

\exe{, difficulty=3}{
	On dit d'une série qu'elle est symétrique autour de 0 si chaque valeur $x$ de la série peut être associée à une unique valeur $-x$ de la série.
	Ainsi $(-1 ; 0 ; 2 ; -2 ; 1)$ est symétrique autour de 0 mais $(1 ; 0 ; -1 ; -1)$ ne l'est pas.
	\begin{enumerate}
		\item
		Montrer qu'une série symétrique autour de 0 a une moyenne nulle.
		\item
		Soit $n\geq1$ un entier naturel.
		À l'aide de la linéarité de la moyenne, montrer que la série statistique $(0 ; 1 ; 2 ; \hdots ; n-1 ; n)$ a pour moyenne $\frac{n}2$.
		\item
		En déduire que
			\[ 1 + 2 + 3 + \dots + (n-1) + n = \dfrac{n(n+1)}2 \]
		\item
		Calculer sans calculatrice la valeur de la somme $1 + 2 + 3 + \hdots+ 9~999 + 10~000$.
	\end{enumerate}
}{exe:somme-triangle}{
	À rendre pour possibilité de points bonus...

	Indication pour la première partie :
%	montrer que la série est symétrique par rapport à $\frac{n}2$.
%	Chaque valeur $\frac{n}2 - k$ a une sœur $\frac{n}2 + k$.
%	Sommer les valeurs sœurs ensemble pour simplifier le calcul de la moyenne.
	calculer la nouvelle moyenne après avoir soustrait $\frac{n}2$ à chaque valeur de la série statistique et utiliser la linéarité de la moyenne.
}

%\exe{,difficulty=2}{
%	Montrer que les solutions $(a;b)$ dont la moyenne statistique est 10 sont les points à coordonnées entières d'une droite.
%	
%	Donner une équation de droite qui n'admet aucun point à coordonnées entières.
%}{exe:appf}{
%
%}

% trop dense en notations
% même si je pense que c'est important de le démontrer pour rassurer son intuition vis-à-vis de la moyenne.
%\exe{,difficulty=2}{
%	Soit $X = (x_1 ; x_2 ; \hdots ; x_N)$ une série statistique dont la valeur minimale est $m$ et la valeur minimale est $M$.
%	Montrer que n'importe quelle moyenne pondérée avec des coefficients positifs prend la forme
%		\[ \overline{X} = f_1 x_1 + f_2 x_2 + \hdots + f_N x_N, \]
%	avec $f_1 + f_2 + \hdots + f_N = 1$ et $0 \leq f_i \leq 1$ pour tout $i$, puis en déduire que
%		\[ m \leq \overline{X} \leq M. \]
%	On dit que la moyenne est une \emph{combinaison convexe} des valeurs.
%}{exe:combinaison-convexe}{
%
%}

%c'est sympa mais bon voilà sans \sum c'est relou non ?
%\exe{,difficulty=2}{
%	Soit $X = (x_1 ; x_2 ; \hdots ; x_n)$ une série statistique.
%	Posons $X^2 = (x_1^2 ; x_2^2 ; \hdots ; x_n^2)$ et $Y = \left((x_1 - \overline{X})^2 ; (x_2 - \overline{X})^2 ; \hdots ; (x_n - \overline{X})^2\right)$.
%	
%	En considérant la moyenne $\overline{Y}$, montrer que $\overline{X^2} \geq \overline{X}^2$.
%	En déduire que $|x|+|y| \leq \sqrt2 \norm{\pvec{x}{y}}$.
%}{exe:variance}{
%
%}

\exe{,difficulty=3}{
	Le but de cet exercice est d'étudier le comportement de l'écart-type d'une série statistique lorsqu'une valeur moyenne est ajoutée à répétition.
	\begin{enumerate}
		\item
		Montrer qu'ajouter une valeur moyenne à une série statistique ne change pas sa moyenne.
		\item
		\begin{multicols}{2}
		%En déduire, sans calcul littéral, la moyenne de la série statistique ci-contre, où $N\in\N$ est un entier naturel.
		Donner la moyenne de la série statistique ci-contre, où $N\in\N$ est un entier naturel.
		
		\begin{center}
		\begin{tabular}{|c|c|c|c|}\hline
		Valeur   & 0 & 20 & 10 \\ \hline
		Effectif & 17 & 17 & $N$ \\ \hline
		\end{tabular}
		\end{center}
		\end{multicols}
	
		\item
		Exprimer l'écart-type $\sigma$ de cette série comme fonction de $N$.
		Calculer $\sigma$ lorsque $N$ vaut 10 ; 20 ; 50 ; 100 à l'aide de la fonction.
		%La série devient-elle plus ou moins concentrée autour de sa moyenne ?
		Vers quelle valeur l'écart-type $\sigma$ s'approche-t-il \mbox{lorsque $N$ grandit ?}
		
		\item
		Donner le premier entier $N\in\N$ à partir duquel l'écart-type est inférieur ou égal à 1.
		\item
		Est-il possible, pour n'importe quel seuil $S>0$, de trouver un entier $N\in\N$ tel que l'écart-type soit inférieur à $S$ ?
	\end{enumerate}
}{exe:variance-exemple}{
	À rendre pour possibilité de points bonus...
}

%%%%%%%%%%%

\newpage
\fancyhead[C]{\textbf{Solutions}}
\shipoutAnswer

\end{document}
